% =====================================================================
%  Estilo de la guía del segundo parcial — Programación 3
%  Tokens de diseño:
%    tinta      #1E2A4A   texto y títulos
%    gris       #5B6375   texto secundario, notas al margen
%    papel      #FFFFFF
%    acento por parte (identifica la parte, no decora):
%      dyc #2F6F8F  pd #7A4E9C  flujo #2E7D5B  np #A33B3B
%  Tipografía: Libertinus Serif (cuerpo + matemática), Libertinus Sans (títulos)
%  Layout: bloque de texto a la izquierda, margen derecho ancho para
%          referencias a parciales/exámenes ("visto en").
% =====================================================================
\documentclass[11pt,a4paper]{article}

\usepackage[T1]{fontenc}
\usepackage[utf8]{inputenc}
\usepackage[spanish,es-noshorthands,es-nodecimaldot]{babel}
\usepackage{libertinus}
\usepackage{amsmath,amssymb,amsthm,mathtools}
\usepackage{libertinust1math}
\usepackage[a4paper,left=2.2cm,right=6.0cm,top=2.4cm,bottom=2.6cm,
            marginparwidth=4.1cm,marginparsep=0.7cm]{geometry}
\usepackage[dvipsnames]{xcolor}
\usepackage{microtype}
\usepackage{enumitem}
\usepackage{booktabs,tabularx,array}
\usepackage{tikz}
\usetikzlibrary{arrows.meta,positioning,shapes,calc,fit,backgrounds}
\usepackage[noend]{algpseudocode}
\usepackage{algorithm}
\usepackage[most]{tcolorbox}
\usepackage{titlesec}
\usepackage{fancyhdr}
\usepackage{marginnote}
\usepackage[hidelinks]{hyperref}
\usepackage{etoolbox}

% ---------- color ----------
\definecolor{tinta}{HTML}{1E2A4A}
\definecolor{gris}{HTML}{5B6375}
\definecolor{grisclaro}{HTML}{E6E8EE}
\definecolor{grafos}{HTML}{9A6A0F}
\definecolor{dyc}{HTML}{2F6F8F}
\definecolor{pd}{HTML}{7A4E9C}
\definecolor{flujo}{HTML}{2E7D5B}
\definecolor{np}{HTML}{A33B3B}
\colorlet{acento}{tinta}           % se redefine al abrir cada parte
\color{tinta}

% ---------- tipografía ----------
\linespread{1.12}
\setlength{\parindent}{0pt}
\setlength{\parskip}{0.45em}
\newcommand{\sans}{\sffamily}

% ---------- títulos ----------
\titleformat{\section}{\sans\Large\bfseries\color{acento}}{\thesection}{0.8em}{}
\titleformat{\subsection}{\sans\large\bfseries\color{tinta}}{\thesubsection}{0.7em}{}
\titleformat{\subsubsection}{\sans\normalsize\bfseries\color{tinta}}{}{0em}{}
\titlespacing*{\section}{0pt}{1.6em}{0.6em}
\titlespacing*{\subsection}{0pt}{1.2em}{0.4em}
\titlespacing*{\subsubsection}{0pt}{0.9em}{0.25em}
\renewcommand{\thesection}{\arabic{section}}

% ---------- encabezado ----------
\pagestyle{fancy}
\fancyhf{}
\fancyheadoffset[R]{4.8cm}
\renewcommand{\headrulewidth}{0pt}
\newcommand{\parteactual}{}
\fancyhead[L]{\sans\footnotesize\color{gris}Programación 3 — guía del segundo parcial}
\fancyhead[R]{\sans\footnotesize\color{acento}\parteactual}
\fancyfoot[R]{\sans\footnotesize\color{gris}\thepage}
\fancyfootoffset[R]{4.8cm}

% ---------- portada de cada parte: la fórmula central es el protagonista ----------
% \parte{nombre-color}{Título}{fórmula (modo matemático display)}{bajada de una o dos líneas}
\newcommand{\parte}[4]{%
  \clearpage
  \colorlet{acento}{#1}%
  \renewcommand{\parteactual}{#2}%
  \thispagestyle{empty}%
  \begin{tikzpicture}[remember picture,overlay]
    \fill[#1] (current page.north west) rectangle ([xshift=0.9cm]current page.south west);
  \end{tikzpicture}%
  \vspace*{3.2cm}
  {\sans\Huge\bfseries\color{#1}#2\par}
  \vspace{2.2cm}
  \noindent\makebox[\textwidth][l]{\begin{minipage}{\dimexpr\textwidth+4.8cm\relax}
    \color{tinta}\fontsize{26}{34}\selectfont
    \[ #3 \]
  \end{minipage}}
  \vspace{1.8cm}

  \begin{minipage}{0.9\textwidth}\large\linespread{1.3}\selectfont\color{gris}#4\end{minipage}\par
  \addcontentsline{toc}{part}{#2}%
  \clearpage
}

% ---------- recuadros con significado ----------
% receta: cómo se redacta la respuesta en el parcial
\newtcolorbox{receta}[1][]{enhanced,breakable,sharp corners,boxrule=0pt,
  frame hidden,colback=acento!6,borderline west={3pt}{0pt}{acento},
  left=10pt,right=8pt,top=6pt,bottom=6pt,
  title={\sans\bfseries Receta para el parcial\ifstrempty{#1}{}{: #1}},
  coltitle=acento,colbacktitle=acento!6,attach title to upper={\par\smallskip}}
% trampa: errores que se corrigen siempre
\newtcolorbox{trampa}[1][]{enhanced,breakable,sharp corners,boxrule=0pt,
  frame hidden,colback=white,borderline west={1.2pt}{0pt}{np},
  left=10pt,right=8pt,top=4pt,bottom=4pt,
  title={\sans\bfseries\color{np}Errores que restan puntos\ifstrempty{#1}{}{: #1}},
  coltitle=np,colbacktitle=white,attach title to upper={\par\smallskip}}
% problema: enunciado de un problema clásico o de examen
\newtcolorbox{problema}[1]{enhanced,breakable,sharp corners,boxrule=0.4pt,
  colframe=grisclaro,colback=white,left=8pt,right=8pt,top=5pt,bottom=5pt,
  fontupper=\small,title={\sans\bfseries #1},coltitle=tinta,colbacktitle=grisclaro}
% resultado: enunciado clave del teórico (lo que se cita "por (x.y)")
\newtcolorbox{resultado}[1][]{enhanced,breakable,sharp corners,boxrule=0pt,frame hidden,
  colback=grisclaro!45,left=10pt,right=8pt,top=6pt,bottom=6pt,
  before upper={\ifstrempty{#1}{}{{\sans\bfseries #1}\par\smallskip}}}

% ---------- notas al margen: dónde se preguntó ----------
% \visto{2P 2019 Ej. 2: ladrón de casas contiguas}
\newcommand{\visto}[1]{\marginnote{\raggedright\sans\footnotesize\color{gris}%
  {\color{acento}\rule[0.15em]{0.9em}{1.4pt}}\hspace{0.35em}#1}}
% nota al margen libre
\newcommand{\margen}[1]{\marginnote{\raggedright\footnotesize\color{gris}#1}}

% ---------- teoremas ----------
\newtheoremstyle{guia}{0.6em}{0.6em}{\itshape}{}{\sans\bfseries}{.}{0.5em}{}
\theoremstyle{guia}
\newtheorem{teorema}{Teorema}[section]
\newtheorem{lema}[teorema]{Lema}
\newtheorem{proposicion}[teorema]{Proposición}
\theoremstyle{definition}
\newtheorem{definicion}[teorema]{Definición}
\newtheorem*{ejemplo}{Ejemplo}
\renewcommand{\qedsymbol}{\color{acento}$\blacksquare$}

% ---------- pseudocódigo en español ----------
\floatname{algorithm}{Algoritmo}
\renewcommand{\algorithmicrequire}{\textbf{Entrada:}}
\renewcommand{\algorithmicensure}{\textbf{Salida:}}
\algrenewcommand\algorithmicwhile{\textbf{mientras}}
\algrenewcommand\algorithmicdo{\textbf{hacer}}
\algrenewcommand\algorithmicfor{\textbf{para}}
\algrenewcommand\algorithmicforall{\textbf{para todo}}
\algrenewcommand\algorithmicif{\textbf{si}}
\algrenewcommand\algorithmicthen{\textbf{entonces}}
\algrenewcommand\algorithmicelse{\textbf{sino}}
\algrenewcommand\algorithmicreturn{\textbf{retornar}}
\algrenewcommand\algorithmicfunction{\textbf{función}}
\algrenewcommand\algorithmicprocedure{\textbf{procedimiento}}

% ---------- macros ----------
\newcommand{\Oh}{\mathcal{O}}
\newcommand{\OPT}{\mathrm{OPT}}
\newcommand{\NP}{\mathsf{NP}}
\newcommand{\Pclass}{\mathsf{P}}
\newcommand{\coNP}{\mathsf{co}\text{-}\mathsf{NP}}
\newcommand{\redp}{\le_{p}}
\newcommand{\prob}[1]{\textsc{#1}}

\setlist{itemsep=0.15em,topsep=0.3em,leftmargin=1.4em}
\setcounter{tocdepth}{1}
